\section{Simulation of Gaussian random variables}

If you want to simulate \(X\sim \normal(0, \Sigma)\) the way to do this
is to find a matrix \(A\) such that \(A A^\transpose = \Sigma\)
and then define
\[
	X = A Y
\]
with \(Y_i\overset{\iid}\sim \normal(0,1)\). This is because
\[
	\Cov(X)
	= \E[XX^\transpose]
	= A \E[YY^\transpose] A^\transpose
	= A \identity A^\transpose
	= \Sigma.
\]
Observe that there are multiple options \(A\) such that \(AA^\transpose = \Sigma\)
unlike the square root for natural numbers!

\begin{enumerate}
	\item The \textbf{cholesky decomposition} uses a lower triangular matrix \(L\)
	such that \(LL^\transpose = \Sigma\). Assumption: \(\Sigma\) is symmetric, positive definite
	\item the \textbf{eigenvalue decomposition} finds an orthonormal matrix \(U\) and 
	diagonal matrix \(D\) such that \(U D U^\transpose = \Sigma\). Assumption: \(\Sigma\) is symmetric.
	Consequently,
	\[
		A = U \sqrt{D}	
	\]
	does the job, where the square root should be viewed as entry-wise. For
	taking the square root we also need positive definiteness (all eigenvalues greater zero).
	\item The \textbf{singular value decomposition} makes no assumptions about \(\Sigma\)
	and finds two orthogonal matrices \(U\) and \(V\) and a diagonal \(D\)
	such that
	\(\Sigma = U D V\).
	If \(\Sigma\) is symmetric, then \(V=U\), however due to numerical inaccuracies
	this is not necessarily the case. The interesting thing is that for simulation
	we can just throw away \(V\) and use \(A= U \sqrt{D}\) again.
\end{enumerate}

If you check the \verb|method| parameter of the
\href{https://numpy.org/doc/stable/reference/random/generated/numpy.random.Generator.multivariate_normal.html#numpy.random.Generator.multivariate_normal}{numpy implementation of the \texttt{multivariate\_normal} distribution} these three
options are exactly what numpy implements.

The cholesky is the fastest, but numerically most unstable variant. The SVD is the slowest
but numerically most stable variant.
Unfortunately, for incremental simulation we are \textbf{forced} to use the cholesky decomposition!